\section{A continuation criterion}\label{sec:continuation}
The notation in this section and the next three is independent of the induction parameters.
Fix a nonconstant Boolean \(f\), use output-sign symmetry to take \(m=\E f\ge0\), and put
\begin{equation}\label{eq:continuation-notation}
\begin{gathered}
s=\sqrt{1-\rho^2},\qquad r=\rho^2=1-s^2,\qquad d=T_\rho f,\\
h=H(m),\qquad b=\E H(d),\qquad v=1-s,\qquad K=\frac rs.
\end{gathered}
\end{equation}
We work with \(.84\le s<1\), so \(0<r\le .2944\).

Define
\begin{equation}\label{eq:energy-definition}
\mathcal E(d)=\E\sum_i D_i d\,D_i\!\left(\frac{d}{H(d)}\right).
\end{equation}
For the parameter \(\tau=-\log\rho\), one has
\(\partial_\tau d=-\mathcal N d\). Since \(H'(z)=-z/H(z)\),
Equation~\eqref{eq:dirichlet} yields
\begin{equation}\label{eq:energy-derivative}
\frac{db}{d\tau}=\mathcal E(d),
\qquad
\frac{ds}{d\tau}=\frac rs=K.
\end{equation}

This identity is the Hellinger edge derivative of
Chen, Gohari, and Nair~\cite[Section~III, Lemma~6]{HellingerDifferential}
in our generator normalization: their parameter has $\rho=e^{-2t}$,
so $\tau=2t$ here.

\begin{lemma}[Crossing criterion]\label{lem:crossing}
Suppose Equation~\eqref{eq:main} is known for \(s\le .84\).
If it fails for some \(s>.84\), then for the same function there is a parameter
\(.84\le s<1\) for which
\begin{equation}\label{eq:crossing}
b=h-1+s,\qquad \mathcal E(d)\le K.
\end{equation}
It is enough to exclude Equation~\eqref{eq:crossing} for functions other than signed coordinates.
\end{lemma}
\begin{proof}
The desired inequality is \(b-h+1-s\ge0\). If this expression is negative at some finite semigroup time, take the left endpoint \(\tau_0\) of the connected component of strict negativity containing that time, restricted to times after \(s=.84\).
Continuity gives equality at \(\tau_0\).
The right difference quotients there are nonpositive, so differentiability and Equation~\eqref{eq:energy-derivative} give the energy inequality in Equation~\eqref{eq:crossing}.
For \(f=\pm x_i\), one has \(m=0\), \(d=\pm\rho x_i\), and \(b=s\), so the desired inequality is an identity for every parameter.
\end{proof}

The following change of variable will be used in the two smaller bias regimes:
\begin{equation}\label{eq:F-definition}
F(x)=\int_0^x(1-z^2)^{-3/4}\,dz,\qquad g=F(d).
\end{equation}
The function \(F\) is odd, increasing, and has finite limits at \(\pm1\).

\begin{lemma}[Basic energy comparison]\label{lem:basic-energy}
With the preceding notation,
\begin{equation}\label{eq:basic-energy}
\mathcal E(d)\ge\sum_i\E(D_i g)^2\ge\Var(g).
\end{equation}
\end{lemma}
\begin{proof}
For the two values \(z_-\le z_+\) of \(d\) on an edge, the edge contributions have the respective forms
\[
\frac14(z_+-z_-)\left(\frac{z_+}{H(z_+)}-\frac{z_-}{H(z_-)}\right),
\qquad
\frac14(F(z_+)-F(z_-))^2.
\]
The derivative of \(x/H(x)\) is \(H(x)^{-3}=(F'(x))^2\), so Cauchy--Schwarz on \([z_-,z_+]\) compares these in the required direction. The case \(z_-=z_+\) follows by continuity. Summing over edges proves the first inequality.
The second follows from Equation~\eqref{eq:dirichlet}, since every nonconstant Fourier degree is at least one.
\end{proof}

The two smaller bias ranges will combine these energy comparisons with
upper bounds on the variance of $d$ and its odd part.  Booleanity supplies
the following bounds.

\begin{lemma}[Boolean spectral bounds]\label{lem:fourier-bounds}
Set
\[
p=\frac{1-m}{2},\qquad
w=\min\left\{2p,\;4\sqrt{\frac32}\,p^{3/2}\right\}.
\]
Then
\begin{align}
W_1(f)&\le w,\label{eq:first-level-bound}\\
\Var(d)&\le r[(1-r)w+rh^2].\label{eq:variance-bound}
\end{align}
If \(d_{\rm o}(x)=(d(x)-d(-x))/2\) is the odd part of \(d\), then
\begin{equation}\label{eq:odd-variance-bound}
\E d_{\rm o}^2\le r[(1-r^2)w+r^2h^2].
\end{equation}
\end{lemma}
\begin{proof}
Write \(L_+=\max(L,0)\) and \(L_-=\max(-L,0)\).
For a homogeneous linear form \(L=\sum_i a_i x_i\) with \(\E L^2=1\), symmetry gives
\(\E L_+^2=1/2\), while expansion gives
\[
\E L^4=3\left(\sum_i a_i^2\right)^2-2\sum_i a_i^4\le3,
\qquad \E L_+^4\le3/2.
\]
Let \(A=\{f=-1\}\), of measure \(p\).
Since \(\E L=0\),
\(\E fL=-2\E[1_A L]\).
The absolute value of the last integral is at most the larger of
\(\E[1_A L_+]\) and \(\E[1_A L_-]\).
Hölder's inequality therefore gives
\[
|\E fL|\le
2\min\left\{\sqrt{p/2},\;(3/2)^{1/4}p^{3/4}\right\}.
\]
Maximizing over \(L\) and squaring proves Equation~\eqref{eq:first-level-bound}.
Parseval and Booleanity give
\(\sum_{k\ge1}W_k(f)=1-m^2=h^2\).
Thus
\[
\Var(d)=\sum_{k\ge1}r^kW_k(f)
\le rW_1(f)+r^2(h^2-W_1(f)),
\]
which proves Equation~\eqref{eq:variance-bound}.
The odd part retains only odd degrees. Every retained degree above one is at least three, so
\[
\E d_{\rm o}^2
\le rW_1(f)+r^3(h^2-W_1(f)),
\]
proving Equation~\eqref{eq:odd-variance-bound}.
\end{proof}

At a hypothetical crossing, the fixed mean lies either in the central band
\(0\le m\le r^2/2\), in the intermediate band \(r^2/2\le m\le .74\), or in the large-bias band \(.74\le m<1\).
The next sections exclude the crossing in the first two bands and prove the target directly in the third.
